% Generated draft: requires mathematical and source review.
\documentclass[11pt]{article}
\usepackage[margin=25mm]{geometry}
\usepackage{fontspec}
\setmainfont{Latin Modern Roman}
\usepackage{amsmath,amssymb,mathtools,unicode-math}
\setmathfont{Latin Modern Math}
\usepackage{enumitem,needspace,xurl,hyperref}
\setlength{\parindent}{0pt}
\setlength{\parskip}{5pt}
\setlength{\emergencystretch}{3em}
\newcommand{\N}{\mathbb{N}}
\newcommand{\Z}{\mathbb{Z}}
\newcommand{\Q}{\mathbb{Q}}
\newcommand{\R}{\mathbb{R}}
\newcommand{\C}{\mathbb{C}}
\newcommand{\F}{\mathbb{F}}
\newcommand{\A}{\mathbb{A}}
\newcommand{\PP}{\mathbb{P}}
\newcommand{\E}{\mathbb{E}}
\newcommand{\eps}{\varepsilon}
\newcommand{\dd}{\,\mathrm{d}}
\begin{document}
\section*{UnsolvedMath ID: 30003613}
\textbf{Source problem code:} catalogue entry
\subsection{Definitions and mathematical statement}
Let \(X\) be a smooth projective Fano variety, so that its anticanonical divisor \(-K_X\) is ample. The supplied record does not specify the ground field, the range of \(q\), or the meaning of a breaking map. The following therefore records the proposed asymptotic subject to unresolved specifications.

Write \(\operatorname{Mor}(\mathbb P^1,X)\) for the scheme parametrizing morphisms from the projective line to \(X\), over the unspecified ground field. For an irreducible component \(W\), the usual interpretation of dominance is that the evaluation morphism
\[
\operatorname{ev}_W:\mathbb P^1\times W\longrightarrow X,
\qquad (t,f)\longmapsto f(t),
\]
is dominant. A morphism \(f\) has anticanonical degree \(d\) if
\[
\deg f^*\mathcal O_X(-K_X)=d.
\]
The record does not specify whether components and dominance are to be interpreted over the ground field or after extending it to an algebraic closure.

Let \(\mathcal W_d^{\circ}(X)\) denote the components of anticanonical degree \(d\) that are dominant and whose family maps do not rationally factor through breaking maps, as required by the source. This last condition remains undefined in the supplied evidence. If the family map means \(\operatorname{ev}_W\), a rational factorization through a map \(b:Y\to X\) would mean a rational map \(h:\mathbb P^1\times W\dashrightarrow Y\) satisfying \(\operatorname{ev}_W=b\circ h\) as rational maps. The record does not identify the maps \(b\) designated as breaking maps or specify which rational factorizations are to be excluded; this description alone does not define the filter.

Subject to these specifications, set
\[
N^{\circ}(X,q,d)=\sum_{W\in\mathcal W_d^{\circ}(X)}q^{\dim W}.
\]
The question is whether, for every smooth projective Fano variety \(X\) under the intended hypotheses and for \(q\) in its intended range, there exists a constant \(C>0\), independent of \(d\), such that
\[
N^{\circ}(X,q,d)\sim Cq^d d^{\rho(X)-1}\qquad(d\to\infty),
\]
that is,
\[
\lim_{d\to\infty}\frac{N^{\circ}(X,q,d)}{Cq^d d^{\rho(X)-1}}=1.
\]
Here \(\rho(X)\) conventionally denotes the Picard number, the rank of the N\'eron--Severi group; the record does not specify whether the arithmetic or geometric Picard number is intended. It also does not specify whether the limit runs through all positive integer degrees or an admissible subset, whether \(q\) is fixed, or the permitted dependence of \(C\) on other parameters. These qualifications must be resolved before this becomes a fully specified conjecture.
\subsection{Short English statement}
The proposed asymptotic counts dominant families of parametrized rational curves on a smooth projective Fano variety, weighting each component by \(q\) to its dimension and excluding families that factor through breaking maps. It predicts growth \(Cq^d d^{\rho(X)-1}\). The supplied record leaves the exclusion rule, ground field, parameter range, and degree convention unresolved.
\subsection{Sources}
\begin{itemize}
\item Supplied local UnsolvedMath registry record, permanent ID \(30003613\), source problem code \texttt{catalogue entry}, title ``Manin-Type Asymptotics for Rational-Curve Components,'' in the Oberwolfach Reports Open Problems collection. The record gives proposed year 2017 and names no individual proposer.
\item Supplied source citation: \emph{Algebraic Geometry: Birational Classification, Derived Categories, and Moduli Spaces} (2018), Oberwolfach Reports, pp.\ 2703--2767, DOI \url{https://doi.org/10.4171/owr/2017/45}. The background identifies Oberwolfach Reports as reports of mathematical workshops and their open-problem discussions.
\item Supplied literature source, labeled ``Oberwolfach Report 45 2017'': \url{https://doi.org/10.4171/owr/2017/45}.
\item Supplied literature source, labeled ``Lehmann Tanimoto 2019'': \url{https://doi.org/10.1112/S0010437X19007103}.
\item Supplied literature source, labeled ``Beheshti Lehmann Riedl Tanimoto 2022'': \url{https://doi.org/10.1016/j.aim.2022.108557}.
\item Literature assessment dated exactly \texttt{2026-08-21}: the supplied curation labels the problem ``partially solved'' and reports substantial upper bounds and special-family results for the conjectural Manin asymptotic after removing breaking maps, without a general establishment. This is supplied dated curation derived from the OWR material, not an independently verified literature assessment. The cited claims and time-sensitive status require independent verification.
\item The registry labels the statement \texttt{corrected\_verified} and says that notation, typography, or omitted definitions were repaired while preserving mathematical content. This is a supplied editorial qualification, not mathematical verification of this draft or resolution of the ambiguities above.
\end{itemize}
\end{document}
